% from ega0-18-fr.tex
\oldpage[0\textsubscript{IV-1}]{53}
est alors \emph{scindée}~; autrement dit, on peut identifier le $A$-bimodule
$E$ à $B\times\mathfrak{J}$, et la multiplication dans $E$ est alors donnée
par
\[
  (b,z)(b',z')=(bb',bz'+zb'+zz'),
\]
$\mathfrak{J}$ étant considéré comme $B$-bimodule au moyen de $s:B\to E$.
\end{env}

\begin{env}[18.1.5]
\label{0.18.1.5-fr}
Avec les notations de (18.1.2), supposons que le $A$-homomorphisme $f$ soit
\emph{surjectif}, autrement dit fasse de $E$ un $A$-anneau \emph{augmenté sur
$B$}. Alors il est clair que $p_2:G\to F$ est aussi \emph{surjectif},
autrement dit définit sur $G$ une structure de $A$-anneau \emph{augmenté sur
$F$}, que l'on appelle \emph{l'image réciproque} par $g:F\to B$ de l'anneau
augmenté $E$.
\end{env}

\begin{proposition}[18.1.6]
\label{0.18.1.6-fr}
Pour que l'image réciproque $G=E\times_BF$ par $g:F\to B$ du $A$-anneau
augmenté $E$ soit un $A$-anneau augmenté trivial, il faut et il suffit qu'il
existe un $A$-homomorphisme $u:F\to E$ rendant commutatif le diagramme
\[
  \xymatrix{
    & F\ar[ld]_u\ar[d]^g\\
    E\ar[r]_f & B
  }
\]
\end{proposition}

\begin{proof}
La condition est évidemment nécessaire, en prenant $u=p_1\circ s$, où
$s:F\to G$ est un $A$-homomorphisme inverse à droite de $p_2$. Inversement,
s'il existe un $A$-homomorphisme $u$ vérifiant la condition de l'énoncé,
l'existence de l'inverse à droite $s$ de $p_2$ résulte de la propriété
universelle du produit fibré (18.1.2) appliquée aux $A$-homomorphismes
$u:F\to E$ et $1_F:F\to F$.

Ce résultat entraîne en particulier que si $E$ est un anneau augmenté
\emph{trivial}, il en est de même de toutes ses images réciproques.
\end{proof}

\begin{env}[18.1.7]
\label{0.18.1.7-fr}
Reprenons la situation décrite en (18.1.2) et (18.1.3)~; on a évidemment sur
$\mathfrak{R}$ une structure de $G$-bimodule, provenant de sa structure de
$F$-bimodule et de l'homomorphisme d'anneaux $p_2:G\to F$. Comme $i_2$ est
bijectif on a en outre une injection
$\theta=j\circ i_2^{-1}:\mathfrak{R}\to G$, qui est un homomorphisme de
$G$-bimodules. Nous allons voir inversement que, lorsqu'on suppose en outre
que $f$ et $g$ sont \emph{surjectifs}, autrement dit que $E$ et $F$ sont des
$A$-anneaux \emph{augmentés} sur $B$, la donnée d'un tel homomorphisme
$\theta$ permet de reconstituer l'anneau augmenté $E$ (sur $B$) à partir de
l'anneau augmenté $G$ (sur $F$).

De façon précise, donnons-nous un $A$-anneau $F$ augmenté sur $B$ et un
$A$-anneau $G$ augmenté sur $F$, les idéaux d'augmentation étant notés
$\mathfrak{R}$ et $\mathfrak{J}'$ respectivement~:
\begin{equation}
\label{0.18.1.7.1-fr}
\tag{18.1.7.1}
  \xymatrix@C=2pc@R=1.5pc{
    & & 0\ar[d] & & \\
    & & \mathfrak{R}\ar[ld]^\theta\ar[d] & & \\
    0\ar[r] & \mathfrak{J}'\ar[r] & G\ar[r]_h & F\ar[r]\ar[d]^g & 0\\
    & & & B\ar[d] & \\
    & & & 0 &
  }
\end{equation}

